\section*{Chapter \ref{ch:b2:deep-time-climate} --- Climate over Geological Time}

\begin{solution}{exo:b2:deep-time-climate:1}
(a) $900 + 37100 + 700 + 3 = \qty{38703}{PgC}$, about 65 times the \qty{591}{PgC} of the air.
(b) $\num{3.5e6}\times0.012011 = \qty{42000}{PgC}$. (c) $591/42000 = \qty{1.4}{\percent}$.
\end{solution}

\begin{solution}{exo:b2:deep-time-climate:2}
(a) Inputs $11.3 + 3 + 12.3 + 8.1 + 2.0 = \qty{36.7}{Tmol.yr^{-1}}$; outputs $26 + 16 =
\qty{42}{Tmol.yr^{-1}}$, equal within their uncertainties. (b) $(11.3 + 3)/36.7 = 0.39$: well
under half; most of the input is old sedimentary carbon. (c) $\num{3.5e6}/42 = \qty{8.3e4}{yr}$.
\end{solution}

\begin{solution}{exo:b2:deep-time-climate:3}
(a) \ce{Mg2SiO4 + 4 CO2 + 2 H2O -> 2 Mg^{2+} + 4 HCO3^- + SiO2}. (b) With \ce{2 Mg^{2+} + 4 HCO3^-
-> 2 MgCO3 + 2 CO2 + 2 H2O}, the net reaction is \ce{Mg2SiO4 + 2 CO2 -> 2 MgCO3 + SiO2}: two moles
of carbon dioxide per mole of forsterite. (c) $\qty{1}{PgC} = 10^{15}/12.011 =
\qty{8.33e13}{mol}$ of carbon, which needs \qty{4.16e13}{mol} of forsterite, $4.16\times10^{13}
\times\qty{140.7}{g} = \qty{5.86}{Pg}$, nearly six billion tonnes.
\end{solution}

\begin{solution}{exo:b2:deep-time-climate:4}
Hirnantian glaciation, 445 to \qty{443}{Ma}, icehouse; greatest ice of the late Palaeozoic, 307 to
\qty{299}{Ma}, icehouse; Turonian, 94 to \qty{89}{Ma}, greenhouse (the warmest of the Phanerozoic);
PETM, 55.9 to \qty{55.7}{Ma}, greenhouse; EECO, 53 to \qty{49}{Ma}, greenhouse; onset of the
Antarctic ice sheet, about \qty{34}{Ma}, start of the present icehouse; northern ice sheets,
\qty{2.7}{Ma}, icehouse.
\end{solution}

\begin{solution}{exo:b2:deep-time-climate:5}
(a) $63000/(8.314\times288^2) = \qty{0.0914}{K^{-1}}$, which rounds to 0.09. (b) $74000/(8.314
\times288^2) = \qty{0.107}{K^{-1}}$. (c) $\ln1.2/0.09 = \qty{2.03}{K}$ and $\ln 1.2/0.107 =
\qty{1.70}{K}$. (d) At the steady state the weathering must equal the degassing, which fixes $e^{k
\Delta T}$ and so $\Delta T$; the sensitivity only sets how much carbon dioxide produces that
warming.
\end{solution}

\begin{solution}{exo:b2:deep-time-climate:6}
Here $k_T = 0.0914$. (a) Exact $\exp[7577.5\times(1/288 - 1/293)] = 1.567$; approximation
$e^{0.0914\times5} = 1.579$: \qty{0.8}{\percent} too high. (b) Exact $e^{2.482} = 11.97$;
approximation $e^{2.741} = 15.50$: \qty{30}{\percent} too high. (c) $\ln(\text{ratio}) = (E_{\text
a}/R)[x/T_0^2 - x/(T_0(T_0 + x))] = (E_{\text a}/R)\,x^2/(T_0^2(T_0 + x))$. Setting it to $\ln1.1$:
$x^2 = 1.043\,(288 + x)$, so $x = \qty{17.9}{K}$. Beyond about \qty{18}{K} of warming the
exponential form overstates the weathering by more than a tenth.
\end{solution}

\begin{solution}{exo:b2:deep-time-climate:7}
(a) $1444/276 = 5.23$, and $3.93\log_2 5.23 = \qty{9.4}{W.m^{-2}}$. (b) Almost exactly the
\qty{9}{W.m^{-2}} of the Sun: the two forcings cancel. (c) $(1444 - 276)/419.5 =
\qty{2.8}{ppm}$ per million years. The \qty{3.4}{ppm} comes from a straight line fitted to all the
data, which weights every estimate rather than the two ends of the smoothed curve.
\end{solution}

\begin{solution}{exo:b2:deep-time-climate:8}
(a) From $3000/20000 = \qty{0.15}{PgC.yr^{-1}}$ to $7000/3000 = \qty{2.3}{PgC.yr^{-1}}$, a range
that contains the IPCC's 0.3 to 1.5. (b) $35\times12.011/44.01 = \qty{9.6}{PgC.yr^{-1}}$, 6 to
32 times the PETM's 0.3 to \qty{1.5}{PgC.yr^{-1}}. (c) $3.93\log_2(2000/900) = \qty{4.5}{W.m^{-2}}$,
and $3\log_2 2.22 = \qty{3.5}{\celsius}$: about seventy per cent of the \qty{5}{\celsius}. Other
feedbacks or another source of warming, such as methane, must make up the rest.
\end{solution}

\begin{solution}{exo:b2:deep-time-climate:9}
(a) Now $1.2\,W_0e^{k\Delta T} = V = W_0$, so $\Delta T = -\ln 1.2/0.09 = \qty{-2.03}{K}$.
(b) $278\times2^{-2.03/3} = \qty{174}{ppm}$. (c) No: at any steady state the weathering equals the
degassing, which has not changed. The higher weatherability is offset by a colder climate.
\end{solution}

\begin{solution}{exo:b2:deep-time-climate:10}
(a) $k = 0.119$: $\tau = 591\times0.693/(0.08\times0.119\times3\times0.138) = \qty{104}{kyr}$.
(b) $S = 8$: $\tau = \qty{51.5}{kyr}$. (c) Both make the thermostat faster, since $\tau \propto
1/(kS)$. $E_{\text a}$ is uncertain by a factor of about two, $S$ on these time scales by a
similar factor (3 against 5 to \qty{8}{K}): both matter about equally.
\end{solution}

\begin{solution}{exo:b2:deep-time-climate:11}
(a) $0.08\times5000 = \qty{400}{PgC}$ stays in the air: $C/C_0 = 991/591 = 1.677$, or
\qty{466}{ppm}; $\Delta T = 3\log_2 1.677 = \qty{2.24}{K}$. (b) $W - W_0 = 0.138\,(e^{0.09\times2.24}
- 1) = \qty{0.0308}{PgC.yr^{-1}}$, \qty{22}{\percent} more weathering; $5000/0.0308 =
\qty{162}{kyr}$. (c) $\tau = \qty{137}{kyr}$ and $\tau\ln10 = \qty{316}{kyr}$. (d) The warming grows
only as the logarithm of the carbon dioxide, so a large excess raises the weathering less than the
linear law assumes; the removal starts more slowly.
\end{solution}

\begin{solution}{exo:b2:deep-time-climate:12}
(a) Near $M^*$, $\mathrm d\Delta T/\mathrm dM = Sf/(A_0\ln 2\,(1 + fM^*/A_0)) = Sf/(A_0\ln 2\,
C^*/C_0)$ and $\mathrm dW/\mathrm dM = kW\,\mathrm d\Delta T/\mathrm dM$, with $W = V$ there; the
inverse is $\tau'$. (b) $C^*/C_0 = 2^{2.03/3} = 1.597$: $\tau' = 137.4\times1.597/1.2 =
\qty{183}{kyr}$. (c) Slower: at high carbon dioxide each added petagram of carbon warms less,
because the forcing is logarithmic, so the weathering responds less to it.
\end{solution}

\begin{solution}{pb:b2:deep-time-climate:1}
\textbf{1.} $1.3 + 1.3 + 4.5 + 4.2 = \qty{11.3}{Tmol.yr^{-1}}$, $\qty{0.136}{PgC.yr^{-1}}$.
\textbf{2.} $11.5\times0.012011 = \qty{0.138}{PgC.yr^{-1}}$.
\textbf{3.} They agree within about \qty{2}{\percent}: silicate weathering balances the volcanic
degassing, as the \omterm{def:b2:deep-time-climate:carbon-cycle}{long-term carbon cycle} requires.
\textbf{4.} $\num{5.1e18}/0.02897 = \qty{1.76e20}{mol}$ of air; $\num{591e15}/12.011 =
\qty{4.92e16}{mol}$ of carbon; $\num{4.92e16}/\num{1.76e20} = \num{2.80e-4}$, \qty{280}{ppm}.
\textbf{5.} $591/0.136 = \qty{4350}{yr}$.
\textbf{6.} $1 + 0.08 = 1.08$: only \qty{8}{\percent} more.
\textbf{7.} $63000/(8.314\times288^2) = \qty{0.0914}{K^{-1}}$, which rounds to 0.09.
\textbf{8.} $e^{0.09} = 1.09$: \qty{9}{\percent} more per kelvin.
\textbf{9.} $\Delta T = S\log_2(1 + fM/A_0)$; its derivative at 0 is $Sf/(A_0\ln2) =
3\times0.08/(591\times0.693) = \qty{5.86e-4}{K.PgC^{-1}}$.
\textbf{10.} $W_0k\times\num{5.86e-4} = 0.138\times0.09\times\num{5.86e-4} = \qty{7.28e-6}{yr^{-1}}$.
\textbf{11.} Its inverse is $A_0\ln2/(fkSW_0)$, the $\tau$ of the proposition, since the
linearised equation is $\mathrm dM/\mathrm dt = -(\mathrm dW/\mathrm dM)\,M$.
\textbf{12.} All four: $\tau$ is inversely proportional to each.
\textbf{13.} $1/\num{7.28e-6} = \qty{1.37e5}{yr}$.
\textbf{14.} $137\times0.09/0.119 = \qty{104}{kyr}$.
\textbf{15.} $k_T = 0.107$: $137\times0.09/0.107 = \qty{115}{kyr}$.
\textbf{16.} Shorter by a factor of 1.5 to 7: the one-box model and its constant $f$ make the
thermostat too fast.
\textbf{17.} $\num{50e6}/\num{1.37e5} = 364$: the thermostat had time to balance the cycle at every
stage, so the decline followed slow changes of the degassing or of the weatherability.
\textbf{18.} $16\times0.137 = \qty{2.2}{ppm}$.
\textbf{19.} $0.08\times5000 = \qty{400}{PgC}$; $278\times991/591 = \qty{466}{ppm}$.
\textbf{20.} $3\log_2 1.677 = \qty{2.24}{K}$.
\textbf{21.} $0.138\,(e^{0.202} - 1) = \qty{0.0308}{PgC.yr^{-1}}$, \qty{22.3}{\percent}.
\textbf{22.} $137\times\ln10 = \qty{316}{kyr}$.
\textbf{23.} The weathering may have been more sensitive (run-off, a larger activation energy, the
slow feedbacks of a larger Earth system sensitivity); and organic carbon on land and at sea may have
regrown and taken carbon out of the air faster than silicate weathering.
\textbf{24.} In this model the \omterm{def:b2:deep-time-climate:thermostat}{silicate-weathering thermostat} removes an excess of carbon with an
e-folding time of about \textbf{\qty[text-series-to-math=true]{137}{} thousand years}.
\end{solution}
